% from ega4-2-fr.tex
\oldpage[IV-2]{138}

\begin{proposition}[6.2.3]
\EGAReaderTarget{ega4.c04.prop.6-2-3}{IV.2.s6.2.3-fr}
Soient $A$, $B$ deux anneaux locaux noethériens, $k$ le corps résiduel de
$A$, $\varphi:A\to B$ un homomorphisme local, $N$ un $B$-module de type fini.
On suppose que $B$ et $N$ sont des $A$-modules plats. Alors on a
\begin{equation}
\dim.\operatorname{proj}_B(N)
=\dim.\operatorname{proj}_{B\otimes_Ak}(N\otimes_Ak).
\tag{6.2.3.1}\label{IV.2.s6.2.3.1-fr}
\end{equation}
\end{proposition}

Considérons en effet une résolution gauche de $N$ par des $B$-modules libres
de type fini
\[
\cdots\to L_i\to L_{i-1}\to\cdots\to L_0\to N\to0.
\]
Comme les $L_i$ et $N$ sont des $A$-modules plats ($B$ étant un $A$-module
plat) il résulte de (2.1.10) que les $Z_i(L_\bullet)$ sont des $A$-modules
plats, que $L_\bullet\otimes_Ak$ est une résolution gauche du
$(B\otimes_Ak)$-module $N\otimes_Ak$ et que l'on a
$Z_i(L_\bullet\otimes_Ak)=Z_i(L_\bullet)\otimes_Ak$ pour tout $i$. Notons que
puisque $B$ est noethérien, les $Z_i(L_\bullet)$ sont des $B$-modules de type
fini, donc les $Z_i(L_\bullet\otimes_Ak)$ sont des $(B\otimes_Ak)$-modules de
type fini. Or, dire qu'un $B$-module (resp. un $(B\otimes_Ak)$-module) de type
fini est plat équivaut à dire qu'il est libre ou encore projectif
$(0_{\mathrm{III}}, 10.1.3)$. Comme les $Z_i(L_\bullet)$ sont des $A$-modules
plats, il résulte d'autre part de $(0_{\mathrm{III}}, 10.2.5)$ (où l'on fait
$C=B$) que, pour que $Z_i(L_\bullet)$ soit un $B$-module plat, il faut et il
suffit que $Z_i(L_\bullet)\otimes_Ak$ soit un $(B\otimes_Ak)$-module plat. Le
plus petit entier $n$ tel que $Z_{n-1}(L_\bullet)$ soit un $B$-module libre est
donc aussi le plus petit entier tel que $Z_{n-1}(L_\bullet\otimes_Ak)$ soit un
$(B\otimes_Ak)$-module libre, ce qui démontre la proposition (0, 17.2.1).

\subsection{Platitude et profondeur}

\begin{proposition}[6.3.1]
\EGAReaderTarget{ega4.c04.prop.6-3-1}{IV.2.s6.3.1-fr}
Soient $A$, $B$ deux anneaux locaux noethériens, $k$ le corps résiduel de
$A$, $\varphi:A\to B$ un homomorphisme local, $M$ un $A$-module de type fini,
$N$ un $B$-module de type fini. Si $N$ est un $A$-module plat $\neq0$, on a
\begin{equation}
\operatorname{prof}_B(M\otimes_AN)=\operatorname{prof}_A(M)
+\operatorname{prof}_{B\otimes_Ak}(N\otimes_Ak).
\tag{6.3.1.1}\label{IV.2.s6.3.1.1-fr}
\end{equation}
\end{proposition}

On peut se borner au cas où $M\neq0$, sans quoi les deux membres de (6.3.1.1)
sont tous deux égaux à $+\infty$. Nous procéderons par récurrence sur l'entier
$n$ égal au second membre de (6.3.1.1) (qui est fini en vertu des hypothèses,
de (0, 16.4.6.2) et du lemme de Nakayama appliqué à $N$). Supposons d'abord
$n=0$; alors, si $\mathfrak m$ et $\mathfrak n$ désignent les idéaux maximaux
de $A$ et $B$ respectivement, $\mathfrak m$ est un idéal premier associé à $M$
et $\mathfrak n/\mathfrak mB$ un idéal premier de
$B/\mathfrak mB=B\otimes_Ak$ associé à $N\otimes_Ak$ (0, 16.4.6, (i)). On
conclut alors de (3.3.1) que $\mathfrak n$ est un idéal premier de $B$ associé
à $M\otimes_AN$, donc (0, 16.4.6, (i)) que
$\operatorname{prof}_B(M\otimes_AN)=0$. Supposons donc $n>0$, et distinguons
deux cas:

$a)$ Supposons $\operatorname{prof}_A(M)>0$. Soit $x\in\mathfrak m$ un
élément $M$-régulier, et posons
\[
A'=A/xA,\qquad B'=B/xB,\qquad M'=M/xM,\qquad N'=N/xN;
\]
on a
\[
B'\otimes_{A'}k=B\otimes_Ak,\qquad N'\otimes_{A'}k=N\otimes_Ak
\]
puisque $x\in\mathfrak m$, et
\[
M'\otimes_{A'}N'=(M\otimes_AN)/x(M\otimes_AN);
\]

